\chapter{API Loading and Error States}

Every component that fetches data from an API passes through the same
handful of states, whether it is a task list, a user profile, or a
dashboard widget. Handling these states consistently is what separates an
app that feels broken during a slow network from one that feels solid.

\section{The State Chain}

At minimum, a data-fetching component moves through \texttt{idle},
\texttt{loading}, and then either \texttt{success} or \texttt{error}. A
well-behaved component never leaves the user staring at a blank screen or a
frozen spinner.

\begin{diagrambox}[API Component State Chain]
\begin{verbatim}
   idle
    |
    |  useEffect fires -> fetch starts
    v
 loading  ------------------------> (spinner / skeleton shown)
    |
    +---------------------+
    |                      |
    v                      v
 success                error
    |                      |
    v                      v
 data.length === 0?   show error message
    |         |            + Retry button
   yes        no                |
    |          |                v
    v          v          user clicks Retry
 empty state  render list  ---> back to loading
\end{verbatim}
\end{diagrambox}

\begin{notebox}[NOTE]
There are really four visible outcomes, not two: \textbf{loading},
\textbf{error}, \textbf{empty} (success but no data), and \textbf{populated}
(success with data). Most bugs in real apps come from forgetting the empty
state and assuming success always means "there is a list to render."
\end{notebox}

\section{Implementing the Pattern}

The same three pieces of state --- data, loading, error --- cover almost
every fetching component in the app.

\begin{lstlisting}[language=JavaScript]
// components/TaskList.jsx
import { useEffect, useState } from "react";
import api from "../api/axios";

function TaskList() {
  const [tasks, setTasks] = useState([]);
  const [loading, setLoading] = useState(true);
  const [error, setError] = useState(null);

  const fetchTasks = async () => {
    setLoading(true);
    setError(null);
    try {
      const res = await api.get("/tasks");
      setTasks(res.data.data);
    } catch (err) {
      setError(err.response?.data?.message || "Failed to load tasks");
    } finally {
      setLoading(false);
    }
  };

  useEffect(() => {
    fetchTasks();
  }, []);

  if (loading) return <Spinner />;

  if (error) {
    return (
      <div className="text-center py-8">
        <p className="text-red-600 mb-3">{error}</p>
        <button onClick={fetchTasks} className="px-4 py-2 bg-blue-600 text-white rounded">
          Retry
        </button>
      </div>
    );
  }

  if (!tasks.length) {
    return <p className="text-gray-500 text-center py-8">No tasks yet. Create your first one!</p>;
  }

  return (
    <ul className="space-y-2">
      {tasks.map((task) => (
        <li key={task._id}>{task.title}</li>
      ))}
    </ul>
  );
}

export default TaskList;
\end{lstlisting}

\begin{warnbox}[WARNING]
Always reset \texttt{error} to \texttt{null} at the start of a new fetch
attempt. Without it, a successful retry can leave a stale error message
rendered alongside (or instead of) the fresh data, because \texttt{error}
never got cleared.
\end{warnbox}

\section{Loading Spinners}

A spinner is the simplest loading indicator --- appropriate for fast,
small requests where showing structure would be overkill.

\begin{lstlisting}[language=JavaScript]
// components/Spinner.jsx
function Spinner() {
  return (
    <div className="flex justify-center items-center py-8">
      <div className="h-8 w-8 border-4 border-blue-600 border-t-transparent
                       rounded-full animate-spin" />
    </div>
  );
}

export default Spinner;
\end{lstlisting}

\section{Skeleton Screens}

For content-heavy views (task cards, profile pages), a skeleton that mimics
the eventual layout reduces perceived loading time better than a spinner,
because the user sees the shape of what's coming.

\begin{lstlisting}[language=JavaScript]
// components/TaskCardSkeleton.jsx
function TaskCardSkeleton() {
  return (
    <div className="border rounded-lg p-4 space-y-3 animate-pulse">
      <div className="h-4 bg-gray-200 rounded w-3/4" />
      <div className="h-3 bg-gray-200 rounded w-1/2" />
      <div className="h-3 bg-gray-200 rounded w-full" />
    </div>
  );
}

// Usage while loading:
// {loading
//   ? Array.from({ length: 4 }).map((_, i) => <TaskCardSkeleton key={i} />)
//   : tasks.map((t) => <TaskCard key={t._id} task={t} />)}
\end{lstlisting}

\begin{tipbox}[TIP]
Use spinners for quick actions (button clicks, form submits) and skeletons
for initial page/list loads. A spinner on a full list load makes the whole
page flash blank-then-full; a skeleton feels smoother because the layout
never jumps.
\end{tipbox}

\section{Choosing the Right State to Show}

\begin{center}
\begin{tabularx}{\textwidth}{lX}
\toprule
\textbf{Situation} & \textbf{What to Render} \\
\midrule
Request in flight & Spinner (short lists/actions) or skeleton (page loads) \\
Request failed & Error message + Retry button, never a blank screen \\
Request succeeded, empty array & Friendly empty state with a call to action (e.g.\ "Create your first task") \\
Request succeeded, has data & The actual list/content \\
\bottomrule
\end{tabularx}
\end{center}

\whatwhyhow
{Every API-driven component tracks \texttt{loading}, \texttt{error}, and \texttt{data} state and renders a different UI for each combination.}
{Without explicit states, slow networks show blank screens, failed requests show nothing or crash the app, and empty results look indistinguishable from a bug.}
{Wrap the fetch in try/catch/finally, set the three state variables accordingly, and branch the render with early returns: loading first, then error (with Retry), then empty, then the populated view.}
